

Cela résulte immédiatement de 1.9 et de la démonstration de 1.5.

\subsection*{2. Complexe cotangent relatif}

\textbf{2.1.} Nous donnerons ici une définition du complexe cotangent qui, sans avoir toute la généralité possible (voir Exp. 0 4.4), sera suffisante pour les besoins du théorème de Riemann--Roch; bien entendu, la définition donnée dans exp. 0 coïncide dans le cas envisagé avec la définition donnée ici.

Soit $f:X\to Y$ un morphisme lissifiable, c'est-à-dire admettant une factorisation :
\[
\begin{tikzcd}[column sep=large,row sep=large]
X \arrow[rr,"i"] \arrow[dr,"f"'] && V \arrow[dl,"g"]\\
& Y, &
\end{tikzcd}
\]
où $g$ est un morphisme lisse, et $i$ une immersion qu'on peut supposer fermée. Soit $\underline{J}$ l'idéal de $X$ dans $V$. La dérivation canonique :
\[
d:\underline{O}_V\longrightarrow\underline{\Omega}^1_{V/Y}
\]
définit un homomorphisme :
\[
\underline{J}\longrightarrow\underline{O}_V\xrightarrow{d}\underline{\Omega}^1_{V/Y}\longrightarrow
\underline{\Omega}^1_{V/Y}\otimes_{\underline{O}_V}\underline{O}_X
\]
qui s'annule sur $\underline{J}^2$, et définit donc un homomorphisme, encore noté $d$ :
\begin{equation}\tag{2.1.1}
\underline{J}/\underline{J}^2\xrightarrow{d} i^*(\underline{\Omega}^1_{V/Y}).
\end{equation}
On vérifie aisément que cet homomorphisme est $\underline{O}_X$-linéaire. On définit alors un complexe de chaînes $\underline{L}^{X/Y}_\bullet$ en posant :
\[
\underline{L}^{X/Y}_0=i^*(\underline{\Omega}^1_{V/Y}),\qquad
\underline{L}^{X/Y}_1=\underline{J}/\underline{J}^2;
\qquad
\underline{L}^{X/Y}_i=0\quad\text{si }i\neq0,1,
\]
et en prenant comme homomorphisme $\underline{L}^{X/Y}_1\to \underline{L}^{X/Y}_0$ l'homomorphisme (2.1.1).

